\documentclass[12pt,a4paper]{article}
\usepackage[utf8]{inputenc}
\usepackage[T1]{fontenc}
\IfFileExists{bahasa.ldf}{\usepackage[bahasa]{babel}}{\renewcommand{\contentsname}{Daftar Isi}\renewcommand{\figurename}{Gambar}\renewcommand{\tablename}{Tabel}\renewcommand{\refname}{Referensi}}
\usepackage[margin=2.5cm]{geometry}
\usepackage{graphicx,float,booktabs,array,tabularx,longtable}
\usepackage{amsmath,amssymb}
\usepackage{xcolor,listings,caption,setspace,hyperref,etoolbox}
\onehalfspacing
\setlength{\parindent}{0pt}\setlength{\parskip}{6pt}
\hypersetup{colorlinks=true,linkcolor=black,urlcolor=blue,citecolor=black}
\captionsetup{font=small,labelfont=bf}
\lstset{
  language=Python,basicstyle=\ttfamily\scriptsize,keywordstyle=\color{blue!70!black},
  commentstyle=\color{green!40!black},stringstyle=\color{orange!70!black},
  numbers=left,numberstyle=\tiny\color{gray},breaklines=true,frame=single,
  showstringspaces=false,tabsize=4,inputencoding=utf8,extendedchars=true,
  literate={→}{{$\rightarrow$}}1 {é}{{e}}1 {–}{{-}}1,
  xleftmargin=1.5em,framexleftmargin=1.5em}
% Gambar aman: bila berkas belum ada, tampil kotak penanda berisi path-nya
\newcommand{\safeimg}[4]{%
  \begin{figure}[H]\centering
  \IfFileExists{#1}{\includegraphics[width=#2\textwidth]{#1}}{%
    \fbox{\parbox{0.85\textwidth}{\centering\vspace{1.2cm}\textbf{[INSERT GAMBAR]}\\ \texttt{\detokenize{#1}}\vspace{1.2cm}}}}
  \caption{#3}\label{#4}\end{figure}}
% Source code aman
\newcommand{\code}[2]{%
  \IfFileExists{#1}{\lstinputlisting[caption={\texttt{\detokenize{#1}} --- #2}]{#1}}{%
    \fbox{\texttt{\detokenize{#1}} belum ditemukan}}}
\newcommand{\todo}[1]{\textcolor{red}{\textbf{[#1]}}}

\begin{document}
\begin{titlepage}
\centering
\begin{figure}
    \centering
    \includegraphics[width=0.5\linewidth]{logo_itb.jpg}
    \label{fig:placeholder}
\end{figure}
\vspace{0.6cm}
{\large\bfseries LAPORAN TUGAS BESAR 2\\ IF3170 INTELIGENSI BUATAN\par}
\vspace{0.4cm}
{\Large\bfseries Implementasi Algoritma Supervised Learning\\ untuk Prediksi Pembatalan Pemesanan Ronaldo Hotel\texttrademark\par}
\vspace{1cm}
{\bfseries Kelompok 25 IMKSIAI}\\[0.6cm]
\begin{tabular}{llr}
\toprule
Nama & NIM \\
\midrule
Yusuf Faishal Listyardi & 13524014  \\
Aziza Dharma Putri & 13524017 \\
Niko Samuel Simanjuntak & 13524029 \\
An-Dafa Anza Alvansyah & 13524038 \\
Kalyca Nathania Benedicta Manullang & 13524071 \\
\bottomrule
\end{tabular}
\vfill
{\bfseries Program Studi Teknik Informatika\\ Sekolah Teknik Elektro dan Informatika\\ Institut Teknologi Bandung\\ November 2026}
\end{titlepage}

\tableofcontents
\newpage

%====================================================================
\section{Deskripsi Persoalan}
\subsection{Latar Belakang dan Tujuan}
Pembatalan pemesanan kamar membuat hotel kehilangan pendapatan karena kamar yang sudah dialokasikan tidak terpakai. Tugas besar ini meminta kelompok membangun model \textit{supervised learning} yang memprediksi apakah sebuah pemesanan di Ronaldo Hotel\texttrademark{} akan dibatalkan (\texttt{is\_canceled} = 1) atau tidak (0). Model dinilai di Kaggle dengan metrik \textit{Macro F1}, yaitu rata-rata F1 kelas ``tidak batal'' dan kelas ``batal'':
\begin{equation}
\text{Macro F1}=\tfrac12\left(F1_{0}+F1_{1}\right),\qquad F1_c=\frac{2P_cR_c}{P_c+R_c}.
\end{equation}
Macro F1 dipilih karena kelas tidak seimbang dan kedua kelas sama pentingnya.

Algoritma yang wajib diimplementasikan \textit{from scratch} dengan NumPy adalah K-Nearest Neighbors (KNN), Decision Tree Learning (DTL), Logistic Regression, Support Vector Machine (SVM), dan Naive Bayes. Hasilnya dibandingkan dengan \textit{scikit-learn}. Seluruh model dapat disimpan dan dimuat kembali, dan minimal satu hasil prediksi diunggah ke Kaggle.

\subsection{Pemahaman Dataset}
Data latih \texttt{train.csv} berisi 27.373 baris dan 19 kolom, sedangkan \texttt{test.csv} berisi 6.844 baris dan 18 kolom (tanpa target). Proporsi kelas pada data latih 67,2\% tidak batal dan 32,8\% batal, sehingga data tidak seimbang secara moderat.

\begin{table}[H]\centering\small
\begin{tabularx}{\textwidth}{lX}
\toprule
Kelompok fitur & Kolom \\
\midrule
Identitas & \texttt{id} (tidak dipakai sebagai fitur) \\
Tamu & \texttt{num\_adults}, \texttt{num\_children}, \texttt{is\_returning\_guest}, \texttt{prior\_cancellations}, \texttt{prior\_completed\_bookings} \\
Menginap & \texttt{weekend\_nights}, \texttt{weekday\_nights}, \texttt{room\_category}, \texttt{meal\_package}, \texttt{needs\_parking}, \texttt{special\_requests} \\
Waktu & \texttt{days\_before\_arrival}, \texttt{arrival\_year}, \texttt{arrival\_month}, \texttt{arrival\_day} \\
Harga \& kanal & \texttt{nightly\_rate\_idr}, \texttt{booking\_channel} \\
Target & \texttt{is\_canceled} \\
\bottomrule
\end{tabularx}
\caption{Kelompok kolom pada dataset.}
\end{table}

Hasil eksplorasi pada notebook menunjukkan hal-hal berikut.
\begin{itemize}
\item Delapan kolom memiliki nilai hilang pada data latih, yaitu \texttt{num\_children} (3.035), \texttt{special\_requests} (2.500), \texttt{nightly\_rate\_idr} (2.038), \texttt{meal\_package} (1.760), \texttt{days\_before\_arrival} (1.123), \texttt{needs\_parking} (1.082), \texttt{room\_category} (1.002), dan \texttt{num\_adults} (809). Pola yang sama muncul pada data test.
\item Terdapat 64 baris duplikat (tanpa \texttt{id}), 413 pemesanan berharga 0 (komplimenter), 106 pemesanan dengan 0 tamu dewasa, dan 64 pemesanan dengan 0 malam menginap. Tidak ada tanggal kedatangan yang tidak valid.
\item Fitur kategorikal adalah \texttt{meal\_package} (4 nilai), \texttt{room\_category} (7 nilai), dan \texttt{booking\_channel} (5 nilai).
\item Analisis PCA eksploratif pada fitur terproses menunjukkan 5, 10, 15, dan 20 komponen utama menjelaskan 51,0\%, 78,8\%, 94,3\%, dan 99,0\% variansi. Variansi tersebar cukup merata sehingga reduksi dimensi tidak memberi pengurangan yang besar.
\end{itemize}

%====================================================================
\section{Pembahasan}
\subsection{Struktur Proyek dan Reproduktibilitas}
Kode berada pada folder \texttt{src/} (modul algoritma), notebook pada \texttt{src/}, hasil model pada \texttt{output/models/}, berkas submission pada \texttt{output/}, dan gambar pada \texttt{output/visualization/}. Semua eksperimen memakai \textit{seed} 42 dan 5-fold CV berstrata yang dibuat dari fungsi \texttt{stratified\_kfold\_indices} pada \texttt{metrics.py}. Seluruh metrik (\texttt{confusion\_matrix}, \texttt{f1\_per\_class}, \texttt{macro\_f1}, \texttt{accuracy}, \texttt{misclassification\_cost}) juga diimplementasikan sendiri dengan NumPy.

\subsection{Data Cleaning dan Preprocessing}
\subsubsection{Pembagian data dan pencegahan kebocoran}
Data latih dibagi 80\%/20\% secara berstrata menjadi 21.898 baris \textit{training split} dan 5.475 baris validasi (proporsi batal 0,328 pada keduanya). Pembagian dilakukan \textbf{sebelum} preprocessing. Semua statistik (median, rata-rata, simpangan baku, kategori, batas diskretisasi) dipelajari hanya dari data latih dan diterapkan ke validasi serta test. Pada CV, preprocessing di-\textit{fit} ulang di dalam setiap fold. Data test hanya dipakai untuk menghasilkan prediksi dan tidak digabung dengan data latih.

\subsubsection{Langkah preprocessing dan alasannya}
\begin{table}[H]\centering\small
\begin{tabularx}{\textwidth}{>{\raggedright\arraybackslash}p{3.6cm}X}
\toprule
Langkah & Keterangan dan alasan \\
\midrule
Penghapusan duplikat & 40 baris duplikat dibuang dari data latih (\texttt{drop\_duplicate\_rows}) agar baris yang sama tidak terhitung dua kali, sehingga \textit{training split} menjadi 21.858 baris. Data test tidak dihapus karena setiap \texttt{id} harus tetap diprediksi. \\
Imputasi numerik & Nilai hilang pada fitur numerik diisi median data latih (misalnya \texttt{days\_before\_arrival}=57, \texttt{nightly\_rate\_idr}=1.647.000). Median tahan terhadap pencilan. \\
Imputasi kategorikal & Nilai hilang diberi kategori \texttt{Unknown}. Hilangnya data bisa membawa informasi sehingga tidak dibuang. \\
Penanganan kategori langka & Kategori yang jarang pada data latih (\texttt{Full Board}, \texttt{Type E}) digabung ke \texttt{Other}. Kategori yang belum pernah terlihat pada test juga masuk \texttt{Other}. \\
Fitur turunan & \texttt{log\_lead\_time} (log dari \texttt{days\_before\_arrival}) untuk meredam kemiringan distribusi, dan \texttt{arrival\_date\_ordinal} dari tahun-bulan-hari agar urutan waktu tercermin dalam satu angka. \\
One-hot encoding & Fitur kategorikal diubah menjadi kolom biner untuk model berbasis jarak dan model linear, sehingga total menjadi 35 fitur pada Logistic Regression. \\
Standardisasi z-score & Dipakai pada KNN, Logistic Regression, dan SVM yang sensitif terhadap skala. DTL memakai 17 fitur asli tanpa one-hot karena menangani fitur campuran dan nilai hilang secara langsung. \\
Diskretisasi kuantil & Khusus Naive Bayes, fitur numerik diubah menjadi \textit{bin} berbasis kuantil sehingga Naive Bayes kategorikal dapat dipakai (sesuai slide Bayesian Network IF3170 hlm.~27--42). \\
\bottomrule
\end{tabularx}
\caption{Ringkasan preprocessing.}
\end{table}
Sesuai hasil QnA dengan asisten dosen, preprocessing dibuat sesama mungkin antarmodel. Perbedaan hanya diberikan bila ada alasan teknis (DTL tidak memerlukan one-hot atau imputasi, Naive Bayes memerlukan diskretisasi, SVM memakai varian \texttt{SVMPreprocessor}), dan keputusan dipilih berdasarkan hasil eksperimen CV, bukan asumsi. Kelas \texttt{Preprocessor}, \texttt{NaiveBayesPreprocessor}, dan \texttt{SVMPreprocessor} masing-masing memiliki \texttt{fit}, \texttt{transform}, \texttt{fit\_transform}, \texttt{save}, dan \texttt{load}.

\code{src/preprocessing.py}{kelas preprocessing}

%--------------------------------------------------------------------
\subsection{K-Nearest Neighbors (KNN)}
\subsubsection{Implementasi}
Kelas \texttt{KNN} pada \texttt{knn.py} menyimpan data latih pada \texttt{fit}, menghitung jarak dengan \texttt{kneighbors}, lalu \texttt{predict\_proba} dan \texttt{predict} melakukan pemungutan suara. Jarak yang dipakai adalah Minkowski dengan $p$ yang dapat diatur (pilihan terbaik $p=1$, yaitu Manhattan). Pada pembobotan jarak, tetangga ke-$i$ berbobot $1/d_i^{\text{power}}$. Metode \texttt{save} dan \texttt{load} menyimpan model.

Karena jarak dipengaruhi semua fitur secara sama, kelompok menambahkan \textit{feature weighting} sebagai metode \textit{wrapper}. Kelas \texttt{FeatureWeightSearch} (\texttt{knn\_weight\_search.py}) memakai pencarian koordinat: bobot satu kelompok fitur diubah satu per satu, dan perubahan dipertahankan bila Macro F1 CV naik.

\subsubsection{Hasil}
Tuning hiperparameter dengan 5-fold CV menghasilkan konfigurasi terbaik $k=31$, $p=1$, bobot jarak, \textit{power}=2 dengan Macro F1 CV 0,8465 ($\pm$0,0049). Bobot fitur memberi selisih rata-rata +0,0088 dan unggul pada 5 dari 5 fold. Setelah pencarian bobot, CV naik dari 0,8465 menjadi 0,8494. Pada validasi Macro F1 naik dari 0,8501 menjadi 0,8541 dan akurasi 0,8732. Bobot akhir terbesar adalah 1,5 pada \texttt{days\_before\_arrival}, \texttt{nightly\_rate\_idr}, \texttt{special\_requests}, dan \texttt{log\_lead\_time}, sedangkan \texttt{num\_children} dibobot 0 dan \texttt{meal\_package}/\texttt{room\_category} dibobot 0,5. Waktu prediksi validasi KNN tergolong lama (104,1 detik) karena semua jarak dihitung langsung. Perbandingan dengan \texttt{KNeighborsClassifier} scikit-learn pada konfigurasi yang sama menghasilkan Macro F1 dan akurasi sama (0,8541 dan 0,8732) dengan kesamaan prediksi 100\%.

\code{src/knn.py}{implementasi KNN}
\code{src/knn_weight_search.py}{pencarian bobot fitur}

%--------------------------------------------------------------------
\subsection{Decision Tree Learning (DTL)}
\subsubsection{Implementasi}
Kelas \texttt{DecisionTree} (\texttt{dtl.py}) mengimplementasikan C4.5. Pemilihan atribut memakai \textit{gain ratio}:
\begin{equation}
\text{Gain}(S,A)=H(S)-\sum_v\frac{|S_v|}{|S|}H(S_v),\quad
\text{GainRatio}=\frac{\text{Gain}(S,A)}{-\sum_v\frac{|S_v|}{|S|}\log_2\frac{|S_v|}{|S|}}.
\end{equation}
Atribut numerik dibelah dengan ambang biner terbaik (dibatasi \texttt{max\_thresholds} kandidat dan diberi penalti Quinlan), atribut kategorikal dibelah per nilai. Nilai hilang ditangani C4.5 dengan menyebarkan sampel secara pecahan ke semua cabang sesuai bobotnya. Pohon dipangkas dengan \textit{cost-complexity pruning} (\texttt{ccp\_alpha}), serta dibatasi \texttt{max\_depth}, \texttt{min\_samples\_split}, \texttt{min\_samples\_leaf}, dan \texttt{min\_gain}. Metode yang tersedia: \texttt{fit}, \texttt{predict\_proba}, \texttt{predict}, \texttt{get\_depth}, \texttt{get\_n\_leaves}, \texttt{save}, \texttt{load}.

\subsubsection{Hasil}
Enam kandidat hiperparameter dievaluasi dengan 5-fold CV. Kandidat terbaik: \texttt{max\_depth}=10, \texttt{min\_samples\_split}=20, \texttt{min\_samples\_leaf}=8, \texttt{min\_gain}=$10^{-5}$, \texttt{ccp\_alpha}=0,0005, dan \texttt{max\_thresholds}=128 (CV 0,8120).
\begin{table}[H]\centering\small
\begin{tabular}{lrrrrr}
\toprule
Model & Macro F1 & Akurasi & Kedalaman & Daun & Waktu fit (s)\\
\midrule
Baseline & 0,79661 & 0,8274 & 6 & 79 & 0,25\\
Tuning CV & 0,81194 & 0,8389 & 10 & 102 & 0,63\\
\bottomrule
\end{tabular}
\caption{Hasil DTL pada data validasi.}
\end{table}
Tuning menaikkan Macro F1 validasi sebesar +0,01533. Fitur dengan kontribusi terbesar adalah \texttt{days\_before\_arrival} (0,404), \texttt{special\_requests} (0,154), \texttt{booking\_channel} (0,152), \texttt{arrival\_month} (0,071), dan \texttt{nightly\_rate\_idr} (0,068). Pembanding \texttt{DecisionTreeClassifier} (kriteria entropi) mencapai Macro F1 0,82569 dan akurasi 0,84895 (waktu fit 0,076 detik) dengan kesamaan prediksi 91,32\%. Selisih Macro F1 \texttt{scratch}$-$\texttt{sklearn} adalah $-0{,}01375$. Selisih ini wajar karena kedua implementasi berbeda (gain ratio dan penanganan nilai hilang pecahan pada C4.5 buatan sendiri, versus kriteria entropi dan \textit{split} CART pada scikit-learn), sehingga struktur pohon tidak identik. Karena scikit-learn tidak mendukung C4.5 murni, perbandingan ini adalah perbandingan terdekat yang tersedia.

\safeimg{dtl_top10.png}{0.95}{Visualisasi pohon keputusan C4.5 hasil tuning (\texttt{top\_n}=10).}{fig:dtl10}
\safeimg{dtl_top5.png}{0.8}{Visualisasi pohon keputusan C4.5 yang disederhanakan (\texttt{top\_n}=5).}{fig:dtl5}

\code{src/dtl.py}{implementasi C4.5 dan Genetic Algorithm}
\code{src/dtl_visualization.py}{visualisasi pohon}

%--------------------------------------------------------------------
\subsection{Logistic Regression}
\subsubsection{Implementasi}
Kelas \texttt{LogisticRegression} (\texttt{logistic\_regression.py}) memodelkan $p(y=1|x)=\sigma(\theta_0+\theta^\top x)$ dengan $\sigma(z)=1/(1+e^{-z})$. Training memakai \textit{stochastic gradient ascent} per sampel pada log-likelihood: $\theta\leftarrow\theta+\eta\,(y_i-p_i)\,x_i$. Regularisasi yang didukung adalah L1 (subgradien), L2, dan Elastic Net:
\begin{equation}
\mathcal{L}(\theta)=-\tfrac1n\sum_i\big[y_i\log p_i+(1-y_i)\log(1-p_i)\big]+\lambda\Big[\alpha\|\theta\|_1+\tfrac{1-\alpha}{2}\|\theta\|_2^2\Big].
\end{equation}
Hyperparameter adalah \texttt{learning\_rate}, \texttt{epochs}, \texttt{penalty}, \texttt{reg\_strength}, \texttt{l1\_ratio}, dan \texttt{optimizer} (\texttt{sga} atau \texttt{adam}). Intersep tidak diregularisasi. Karena L1 diterapkan dengan subgradien, koefisien jarang tepat bernilai nol sehingga efek seleksi fitur tidak diklaim.

\subsubsection{Hasil}
\begin{table}[H]\centering\small
\begin{tabular}{llrrr}
\toprule
Penalti & Implementasi & Macro F1 & Akurasi & Log-loss\\
\midrule
none & scratch & 0,76233 & 0,79909 & 0,43467\\
none & scikit-learn & 0,76343 & 0,79890 & 0,43402\\
l1 & scratch & 0,75038 & 0,79598 & 0,44537\\
l1 & scikit-learn & 0,75406 & 0,79689 & 0,44535\\
l2 & scratch & 0,75839 & 0,79872 & 0,43792\\
l2 & scikit-learn & 0,76001 & 0,79890 & 0,43749\\
elasticnet & scratch & 0,75404 & 0,79744 & 0,44130\\
elasticnet & scikit-learn & 0,75611 & 0,79744 & 0,44114\\
\bottomrule
\end{tabular}
\caption{Perbandingan Logistic Regression scratch dan scikit-learn pada validasi.}
\end{table}
Tuning 14 kandidat dengan 5-fold CV memilih \texttt{learning\_rate}=0,0005, \texttt{epochs}=30, penalti \texttt{none} (Macro F1 CV 0,76362). Pada validasi, tuning parameter memberi +0,00147 (0,76381) dan ambang OOF 0,45 memberi $-0{,}00032$ (0,76349). Karena itu ambang 0,5 dipakai untuk model ini. Hasil scratch dan scikit-learn berselisih tipis (sekitar 0,001--0,004 Macro F1), kemungkinan karena penyelesaian optimasi yang berbeda (SGA per sampel dengan jumlah epoch terbatas vs. solver scikit-learn). Pengaruh regularisasi tidak menguntungkan pada konfigurasi yang diuji.

\code{src/logistic_regression.py}{implementasi Logistic Regression}

%--------------------------------------------------------------------
\subsection{Support Vector Machine (SVM)}
\subsubsection{Implementasi}
Kelas \texttt{SVM} (\texttt{svm.py}) menyelesaikan masalah dual soft-margin
\begin{equation}
\max_\alpha\ \sum_i\alpha_i-\tfrac12\sum_{i,j}\alpha_i\alpha_jy_iy_jK(x_i,x_j),\quad 0\le\alpha_i\le C,\ \sum_i\alpha_iy_i=0
\end{equation}
dengan algoritma \textit{Sequential Minimal Optimization} (SMO) dan pemilihan pasangan variabel \textit{working set selection} (Fan dkk.). Fungsi keputusan $f(x)=\sum_i\alpha_iy_iK(x_i,x)+b$. Kernel yang tersedia:
\begin{itemize}
\item Linear: $K(x,z)=x^\top z$
\item Polynomial: $K(x,z)=(\gamma x^\top z+r)^d$
\item Gaussian (RBF): $K(x,z)=\exp(-\gamma\|x-z\|^2)$
\end{itemize}
Sesuai QnA, solver dual boleh memakai pustaka, tetapi kelompok mengimplementasikan SMO sendiri dengan NumPy.

\subsubsection{Hasil}
\begin{table}[H]\centering\small
\begin{tabular}{lrrrrr}
\toprule
Kernel & Macro F1 scratch & Macro F1 sklearn & Akurasi & Kesamaan prediksi & Waktu fit scratch (s)\\
\midrule
Linear & 0,7599 & 0,7599 & 0,7978 & 100\% & 150,0\\
Polynomial & 0,7771 & 0,7771 & 0,8137 & 100\% & 17,6\\
Gaussian & 0,7933 & 0,7931 & 0,8267 & 99,98\% & 62,8\\
\bottomrule
\end{tabular}
\caption{Perbandingan SVM scratch dan \texttt{SVC} scikit-learn.}
\end{table}
Semua SMO konvergen. Kernel Gaussian terbaik (F1 kelas batal 0,710), sehingga dipakai untuk optimasi lanjutan. Hasil scratch hampir identik dengan \texttt{SVC} (selisih Macro F1 Gaussian hanya 0,0002).

\safeimg{output/visualization/svm/svm_linear.png}{0.7}{Batas keputusan SVM kernel linear (2 fitur, 600 sampel).}{fig:svml}
\safeimg{output/visualization/svm/svm_polynomial.png}{0.7}{Batas keputusan SVM kernel polynomial.}{fig:svmp}
\safeimg{output/visualization/svm/svm_gaussian.png}{0.7}{Batas keputusan SVM kernel Gaussian.}{fig:svmg}
Ketiga gambar memakai dua fitur terstandardisasi (\texttt{days\_before\_arrival} dan \texttt{nightly\_rate\_idr}). Garis penuh adalah $f(x)=0$, garis putus-putus margin $f(x)=\pm1$, dan lingkaran menandai \textit{support vector} (362, 350, dan 348 dari 600 titik untuk linear, polynomial, dan Gaussian). Kernel linear menghasilkan garis lurus, sedangkan polynomial dan Gaussian menghasilkan batas melengkung. Kedua kelas masih bercampur pada dua fitur ini.

\code{src/svm.py}{implementasi SVM dan Differential Evolution}

%--------------------------------------------------------------------
\subsection{Naive Bayes}
\subsubsection{Implementasi}
Kelas \texttt{NaiveBayes} (\texttt{naive\_bayes.py}) adalah Naive Bayes kategorikal dengan aturan MAP:
\begin{equation}
\hat y=\arg\max_c\ \Big[\log P(c)+\sum_{j}\log P(x_j|c)\Big],\quad
P(x_j=v|c)=\frac{N_{c,j,v}+\alpha}{N_c+\alpha\,|V_j|}.
\end{equation}
Smoothing Laplace memakai $\alpha=1$ pada model dasar, sesuai arahan asisten dosen (slide Bayesian Network IF3170 2024 hlm.~27--42). Probabilitas dihitung dalam log untuk menghindari \textit{underflow} dan dinormalisasi dengan \textit{log-sum-exp} pada \texttt{predict\_proba}. Fitur numerik didiskretisasi berbasis kuantil oleh \texttt{NaiveBayesPreprocessor}.

\subsubsection{Hasil}
Tuning \texttt{n\_bins}$\in\{5,10,20,40\}$ dengan $\alpha=1$ memilih 40 \textit{bin} (CV 0,7477 $\pm$ 0,0043). Pada validasi Macro F1 0,7463 dan akurasi 0,7821. Implementasi scratch identik dengan \texttt{CategoricalNB} (prediksi 100\% sama, selisih probabilitas maksimum $4{,}3\times10^{-15}$).
\begin{table}[H]\centering\small
\begin{tabular}{lcc}
\toprule
 & Prediksi 0 & Prediksi 1\\
\midrule
Aktual 0 & 3169 & 512\\
Aktual 1 & 681 & 1113\\
\bottomrule
\end{tabular}
\caption{Confusion matrix Naive Bayes dasar pada validasi (F1 per kelas 0,8416 dan 0,6511).}
\end{table}

\code{src/naive_bayes.py}{implementasi Naive Bayes dan PSO}

%--------------------------------------------------------------------
\subsection{Perbandingan Seluruh Model}
\begin{table}[H]\centering\small
\begin{tabular}{lrr}
\toprule
Model (scratch) & Macro F1 validasi & Akurasi\\
\midrule
KNN (bobot fitur) & 0,8541 & 0,8732\\
SVM Gaussian, optimasi lanjutan & 0,8264 & 0,8508\\
DTL C4.5, Genetic Algorithm & 0,8217 & 0,8440\\
DTL C4.5, tuning CV & 0,8119 & 0,8389\\
SVM Gaussian, DE & 0,8242 & 0,8484\\
Logistic Regression (SGA) & 0,7638 & 0,7987\\
Naive Bayes, PSO & 0,7536 & 0,7837\\
Naive Bayes dasar & 0,7463 & 0,7821\\
\bottomrule
\end{tabular}
\caption{Ringkasan Macro F1 pada data validasi 20\% (5.475 baris).}
\end{table}
Pada dataset ini KNN berbobot memberi skor lokal tertinggi, diikuti SVM Gaussian dan DTL. Model linear dan probabilistik (Logistic Regression, Naive Bayes) lebih rendah karena asumsi bentuk yang lebih kaku. Karena KNN, SVM, dan DTL pernah dipakai memilih konfigurasi pada validasi yang sama, skor validasi tersebut adalah estimasi lokal dan bukan estimasi independen. Skor leaderboard Kaggle adalah ukuran akhir.

\subsubsection{Analisis Kesalahan KNN}
Confusion matrix KNN final pada validasi: TN=3383, FP=298, FN=396, TP=1398, dengan F1 per kelas 0,907 dan 0,801. Tingkat kesalahan tertinggi muncul pada \textit{lead time} 90--180 hari (0,179) dan 30--90 hari (0,154), sedangkan kanal \textit{Online Travel Agent} (0,160) dan \textit{Airline} (0,188, hanya 16 sampel) paling sering salah dibanding \textit{Offline Agent} (0,068) dan \textit{Corporate} (0,068).

%====================================================================
\section{Eksperimen Bonus}
\subsection{Analisis Cost-Sensitive (Bonus)}
Pada skenario hotel, prediksi batal yang salah (FP) memiliki biaya 2 dan prediksi tidak batal yang salah (FN) bernilai 1. Matriks biaya:
\begin{center}\small\begin{tabular}{lcc}\toprule & Prediksi 0 & Prediksi 1\\\midrule Aktual 0 & 0 & 2\\ Aktual 1 & 1 & 0\\\bottomrule\end{tabular}\end{center}
Menurut Elkan (2001), ambang optimal adalah $P(\text{batal})>\frac{C_{FP}}{C_{FP}+C_{FN}}=\frac23\approx0{,}667$. Hasil pada validasi (model KNN final):
\begin{table}[H]\centering\small
\begin{tabular}{lrrrrrrr}
\toprule
Aturan & Ambang & FP & FN & Biaya & Biaya/pesanan & Recall batal & Macro F1\\
\midrule
Default & 0,500 & 298 & 396 & 992 & 0,1812 & 0,7793 & 0,8541\\
Bayes 2/3 & 0,667 & 139 & 630 & 908 & 0,1658 & 0,6488 & 0,8269\\
Biaya min. CV & 0,710 & 111 & 700 & 922 & 0,1684 & 0,6098 & 0,8138\\
\bottomrule
\end{tabular}
\caption{Perbandingan aturan ambang berbasis biaya.}
\end{table}
Ambang teoretis 2/3 menurunkan biaya sebesar 7,1\% dibanding ambang default, tetapi menurunkan Macro F1 (0,8541 menjadi 0,8269) karena mengorbankan \textit{recall}. Ambang dengan biaya minimum dari CV adalah 0,71, sedangkan ambang dengan Macro F1 CV tertinggi tetap 0,50. Untuk submission kompetisi yang dinilai Macro F1, ambang 0,5 dipertahankan. Analisis ini hanya relevan bila hotel benar-benar menilai biaya FP dua kali FN.

\subsection{Differential Evolution untuk SVM}
Differential Evolution (Storn \& Price, 1997) diimplementasikan di \texttt{svm.py}. Parameter yang dicari adalah $C$ dan $\gamma$ (skala log) serta ambang keputusan, dengan pencarian pada 4.000 data, tiga fold, delapan kandidat, dan empat generasi. Skor CV terbaik naik dari 0,7665 menjadi 0,7893. CV konfirmasi pada 8.000 data memberi baseline 0,7738 dan kandidat terpilih ($C\approx16{,}02$, $\gamma\approx0{,}06084$) 0,8021. Pada validasi, parameter DE dengan ambang 0 memperoleh Macro F1 0,8242 (baseline 0,7933, +0,0309), sedangkan ambang hasil CV ($-0{,}0755$) memperoleh 0,8211. Ambang 0 dipilih untuk submission. \texttt{SVC} dengan parameter sama memberi 0,8244 dengan kesamaan prediksi 99,98\%.

Optimasi lanjutan menguji representasi fitur (log riwayat, biner, tanggal ordinal dan siklik), pembobotan fitur dan kelas, serta $C$ dan $\gamma$. Representasi tanggal ordinal (0,8002) dan siklik (0,8057) lebih rendah dari baseline (0,8091), dan bobot kelas yang lebih besar menurunkan CV. Pembobotan fitur menghasilkan kandidat terbaik. Pada CV seluruh \textit{training split}, Macro F1 naik dari 0,8211 menjadi 0,8249, dan pada validasi dari 0,8242 menjadi 0,8264 (+0,0022).

\subsection{Particle Swarm Optimization untuk Naive Bayes}
PSO (Kennedy \& Eberhart, 1995; inertia dari Shi \& Eberhart, 1998) diimplementasikan di \texttt{naive\_bayes.py}. Parameter yang dicari adalah jumlah \textit{bin} (5--80), $\log_{10}\alpha$ ($-3$ sampai 1), dan ambang kelas batal (0,2--0,8). Digunakan 12 partikel, 12 iterasi, \textit{seed} 42, inertia 0,7, koefisien personal dan sosial 1,5. PSO memilih 47 \textit{bin}, $\alpha\approx0{,}0567$, ambang $\approx0{,}4804$. Macro F1 CV naik dari 0,7477 menjadi 0,7503, dan Macro F1 validasi dari 0,7463 menjadi 0,7536 (+0,0073). F1 kelas batal naik dari 0,6511 ke 0,6674, FN turun dari 681 ke 606 dan FP naik dari 512 ke 578. Prediksi identik dengan \texttt{CategoricalNB} yang memakai parameter sama. PSO bersifat heuristik dan tidak menjamin optimum global.

\subsection{Genetic Algorithm untuk DTL}
Genetic Algorithm (Holland, 1975; Goldberg, 1989) diimplementasikan di \texttt{dtl.py}. Setiap kandidat berisi kedalaman maksimum, ukuran minimum daun, kekuatan pruning, dan ambang prediksi. Seleksi memakai \textit{tournament}, lalu \textit{crossover}, mutasi, dan \textit{elitism}. Populasi 6 dan 4 generasi dengan fungsi tujuan Macro F1 CV 5-fold. Hasil: \texttt{max\_depth}=10, \texttt{min\_samples\_leaf}=10, \texttt{ccp\_alpha}$\approx$0,000579, ambang $\approx$0,4310. Macro F1 CV naik dari 0,8120 menjadi 0,8186 dan pada validasi 0,8217 (naik +0,0097 dari tuning biasa, F1 kelas batal 0,7585). Peningkatan CV tidak otomatis berarti peningkatan validasi, dan hasil bergantung pada ruang pencarian kecil.

\subsection{Optimizer Adam untuk Logistic Regression}
Adam (Kingma \& Ba, 2015) menyimpan rata-rata bergerak gradien ($m_t$) dan kuadrat gradien ($v_t$) per parameter:
\begin{align}
m_t&=\beta_1m_{t-1}+(1-\beta_1)g_t,\quad v_t=\beta_2v_{t-1}+(1-\beta_2)g_t^2,\\
\hat m_t&=\frac{m_t}{1-\beta_1^t},\quad \hat v_t=\frac{v_t}{1-\beta_2^t},\quad
\theta_t=\theta_{t-1}-\eta\frac{\hat m_t}{\sqrt{\hat v_t}+\epsilon}.
\end{align}
\begin{table}[H]\centering\small
\begin{tabular}{lrrrrrrr}
\toprule
Optimizer & $\eta$ & Epoch & Log-loss latih & Log-loss val & Macro F1 val & Waktu (s) & Epoch ke target\\
\midrule
SGA & 0,0005 & 30 & 0,43288 & 0,43461 & 0,76381 & 3,06 & 3\\
Adam & 0,0010 & 30 & 0,43263 & 0,43453 & 0,76324 & 9,83 & 1\\
\bottomrule
\end{tabular}
\caption{Perbandingan SGA dan Adam (target log-loss latih 0,44). Waktu diukur pada satu eksekusi dan dapat berubah antarperangkat.}
\end{table}
Adam mencapai target log-loss pada epoch pertama, tetapi tiap langkah lebih mahal sehingga waktu total lebih lama. Macro F1 kedua optimizer hampir sama, sehingga Adam bermanfaat pada kecepatan konvergensi per epoch, bukan pada akurasi akhir.

\subsection{Visualisasi Kontur Log-Loss}
Model satu fitur (\texttt{days\_before\_arrival} terstandardisasi) dengan parameter $\theta_0$ dan $\theta_1$ dilatih tanpa regularisasi. Log-loss turun dari 0,69315 pada $(0,0)$ menjadi 0,53957 setelah 32 epoch. Video lintasan parameter tersedia pada \texttt{output/visualization/logistic\_regression/loss\_contour.mp4}. Bobot positif berarti lead time yang lebih panjang menaikkan peluang pembatalan.
\safeimg{output/visualization/logistic_regression/loss_contour_frame.png}{0.9}{Kontur log-loss dan lintasan parameter Logistic Regression (cuplikan dari \texttt{loss\_contour.mp4}).}{fig:contour}

\subsection{Visualisasi Graf Naive Bayes}
Node pusat adalah \texttt{is\_canceled}, panah menuju setiap fitur menggambarkan struktur Naive Bayes dan ketebalannya menunjukkan \textit{mutual information}. Garis putus-putus menandai lima pasangan fitur dengan MI terbesar dan bukan ketergantungan tambahan di dalam model.
\safeimg{output/visualization/naive_bayes/nb_graph.png}{0.95}{Graf Naive Bayes dengan ketebalan panah berdasarkan mutual information.}{fig:nbgraph}

\subsection{Penyimpanan dan Pemuatan Model}
Setiap model memiliki \texttt{save} dan \texttt{load}. Pengujian pada notebook menunjukkan prediksi (dan probabilitas) sebelum dan sesudah dimuat identik untuk KNN, DTL, Logistic Regression, SVM (ketiga kernel), dan Naive Bayes. Model disimpan di \texttt{output/models/}.

\subsection{Submission Kaggle}
Berkas submission dihasilkan di folder \texttt{output/} (6.844 baris, kolom \texttt{id} dan \texttt{is\_canceled}).
\begin{table}[H]\centering\small
\begin{tabular}{lrr}
\toprule
Berkas & Proporsi prediksi batal & Skor Kaggle (Macro F1)\\
\midrule
\texttt{submission\_knn.csv} & 0,299 & \todo{isi}\\
\texttt{submission\_svm\_refined.csv} & \todo{isi} & \todo{isi}\\
\texttt{submission\_dtl\_optimized.csv} & \todo{isi} & \todo{isi}\\
\texttt{submission\_naive\_bayes\_optimized.csv} & \todo{isi} & \todo{isi}\\
\texttt{submission\_logistic\_regression.csv} & \todo{isi} & \todo{isi}\\
\bottomrule
\end{tabular}
\caption{Daftar submission. Isi skor dari halaman \textit{Submissions} Kaggle.}
\end{table}
Nama tim Kaggle: \todo{isi}. Peringkat leaderboard akhir: \todo{isi}.

%====================================================================
\section{Kesimpulan dan Saran}
\subsection{Kesimpulan}
\begin{enumerate}
\item Lima algoritma (KNN, DTL C4.5, Logistic Regression, SVM, Naive Bayes) berhasil diimplementasikan dari awal dengan NumPy, termasuk metrik, CV berstrata, \textit{save/load}, dan pipeline preprocessing tanpa kebocoran data.
\item KNN, SVM, dan Naive Bayes hampir identik dengan scikit-learn (kesamaan prediksi 99,98\%--100\%). Logistic Regression sedikit berbeda (selisih Macro F1 sekitar 0,001--0,004) dan DTL C4.5 berbeda 0,0138 dari \texttt{DecisionTreeClassifier} karena perbedaan algoritma pemisahan.
\item Pada validasi lokal, KNN berbobot terbaik (Macro F1 0,8541), diikuti SVM Gaussian (0,8264) dan DTL dengan Genetic Algorithm (0,8217). Logistic Regression (0,7638) dan Naive Bayes (0,7536) lebih rendah.
\item Fitur yang paling berpengaruh adalah \texttt{days\_before\_arrival}, \texttt{special\_requests}, \texttt{booking\_channel}, dan \texttt{nightly\_rate\_idr}.
\item Optimasi bonus (DE, PSO, GA, Adam) memberi peningkatan kecil sampai sedang pada CV, tetapi tidak selalu pada validasi. Analisis \textit{cost-sensitive} menunjukkan ambang 2/3 menurunkan biaya 7,1\% dengan harga turunnya Macro F1.
\end{enumerate}
\subsection{Saran}
\begin{itemize}
\item Memakai data validasi terpisah (\textit{nested CV}) agar estimasi akhir tidak bias dari pemilihan model.
\item Mempercepat KNN dengan struktur indeks (misalnya KD-Tree) karena prediksi validasi memakan 104 detik.
\item Menambah ruang pencarian dan jumlah generasi/iterasi untuk optimasi GA, PSO, dan DE, serta menguji rekayasa fitur tambahan.
\item Mencoba metode \textit{ensemble} (misalnya \textit{random forest} atau \textit{gradient boosting}) untuk memperbaiki Macro F1.
\end{itemize}

%====================================================================
\section{Pembagian Tugas}
\begin{table}[H]\centering\small
\begin{tabularx}{\textwidth}{>{\raggedright\arraybackslash}p{4cm}lX}
\toprule
Nama & NIM & Kontribusi \\
\midrule
Kalyca Nathania Benedicta Manullang & \todo{NIM} & Genetic Algorithm DTL, \todo{lengkapi} \\
Niko Samuel & \todo{NIM} & \todo{lengkapi} \\
\todo{Anggota 3} & \todo{NIM} & \todo{lengkapi} \\
\todo{Anggota 4} & \todo{NIM} & \todo{lengkapi} \\
\bottomrule
\end{tabularx}
\caption{Pembagian tugas anggota kelompok.}
\end{table}

%====================================================================
\begin{thebibliography}{99}
\bibitem{spec} Tim Asisten IF3170, \textit{Spesifikasi Tugas Besar 2 IF3170 Inteligensi Buatan}, Institut Teknologi Bandung, 2026.
\bibitem{bayes} Slide kuliah IF3170, \textit{Bayesian Network}, hlm.~27--42, 2024.
\bibitem{c45} J. R. Quinlan, \textit{C4.5: Programs for Machine Learning}, Morgan Kaufmann, 1993.
\bibitem{smo} R.-E. Fan, P.-H. Chen, dan C.-J. Lin, ``Working Set Selection Using Second Order Information for Training SVM,'' \textit{JMLR}, vol.~6, 2005.
\bibitem{guide} C.-W. Hsu, C.-C. Chang, dan C.-J. Lin, \textit{A Practical Guide to Support Vector Classification}, National Taiwan University.
\bibitem{elkan} C. Elkan, ``The Foundations of Cost-Sensitive Learning,'' \textit{IJCAI}, 2001.
\bibitem{lipton} Z. C. Lipton, C. Elkan, dan B. Naryanaswamy, ``Optimal Thresholding of Classifiers to Maximize F1 Measure,'' 2014, arXiv:1402.1892.
\bibitem{de} R. Storn dan K. Price, ``Differential Evolution,'' \textit{Journal of Global Optimization}, vol.~11, 1997.
\bibitem{pso} J. Kennedy dan R. Eberhart, ``Particle Swarm Optimization,'' \textit{IEEE ICNN}, 1995.
\bibitem{shi} Y. Shi dan R. Eberhart, ``A Modified Particle Swarm Optimizer,'' \textit{IEEE ICEC}, 1998.
\bibitem{ga} J. H. Holland, \textit{Adaptation in Natural and Artificial Systems}, University of Michigan Press, 1975; D. E. Goldberg, \textit{Genetic Algorithms in Search, Optimization, and Machine Learning}, Addison-Wesley, 1989.
\bibitem{adam} D. P. Kingma dan J. Ba, ``Adam: A Method for Stochastic Optimization,'' 2015, arXiv:1412.6980.
\bibitem{sk} F. Pedregosa dkk., ``Scikit-learn: Machine Learning in Python,'' \textit{JMLR}, vol.~12, 2011.
\end{thebibliography}

\appendix
\section{Lampiran: Kode Program Tambahan}
\code{src/metrics.py}{metrik evaluasi dan CV berstrata}
\code{src/logistic_regression_visualization.py}{kontur log-loss}
\code{src/naive_bayes_visualization.py}{graf Naive Bayes}
\end{document}